\documentclass[11pt]{article}
\input{style-pc}
\usepackage{array}

\begin{document}

\entetesujet{Sujet bac 2020 -- Physique-Chimie -- Série C}

% ============================ CHIMIE ============================
\matiere{CHIMIE}{8 points}

\partie{A : vérification des connaissances}

\noindent\textbf{1. Question à réponse courte.} Soit la réaction chimique équilibrée exothermique suivante :
\[ \ch{2\,SO_{2(g)} + O_{2(g)}} \rightleftharpoons \ch{2\,SO_{3(g)}} \]
Dans quel sens l'équilibre se déplace-t-il :
\begin{enumerate}[label=\alph*.]
  \item lorsqu'on augmente la température ?
  \item lorsqu'on augmente la pression ?
\end{enumerate}

\medskip
\noindent\textbf{2. Texte à trous.} Recopie et complète la phrase ci-après par quatre mots manquants que tu trouveras parmi les six mots suivants : \textit{substance, initialement, demi-vie, moitié, échantillon, totalité}.

\textit{« La $\cdots\cdots$ est la durée au bout de laquelle la $\cdots\cdots$ des noyaux $\cdots\cdots$ présents dans un $\cdots\cdots$ a disparu ».}

\medskip
\noindent\textbf{3. Appariement.} Relie un élément-question de la colonne A à un élément-réponse de la colonne B. Exemple : $a_5 = b_7$.

\begin{center}\renewcommand{\arraystretch}{1.5}
\begin{tabular}{|p{0.38\linewidth}|p{0.42\linewidth}|}
\hline
\textbf{Colonne A} & \textbf{Colonne B} \\
\hline
$a_1$) monobase forte & $b_1$) pH $= -\log 2C$ \\
\hline
$a_2$) diacide fort & $b_2$) pH $= -\log C$ \\
\hline
$a_3$) dibase forte & $b_3$) pH $= 14 + \log 2C$ \\
\hline
$a_4$) mono-acide fort & $b_4$) pH $= 14 + \log C$ \\
\hline
\end{tabular}
\end{center}

\partie{B : application des connaissances}
Les énergies des différents niveaux de l'atome d'hydrogène, exprimées en électrons-volts, sont données par la formule : $E_n = -\dfrac{13{,}6}{n^2}$ où $n$ est un entier naturel non nul.

\begin{enumerate}
  \item \begin{enumerate}[label=\alph*.]
    \item Calcule, en électrons-volts, les énergies correspondant à $n = 1$ ; $n = 2$ ; $n = 3$ et $n = \infty$.
    \item Représente ces quatre niveaux dans un diagramme d'énergie. Échelle : 1 cm pour 1 eV.
  \end{enumerate}
  \item \begin{enumerate}[label=\alph*.]
    \item Calcule l'énergie minimale que l'on doit fournir à un atome d'hydrogène pour qu'il passe de l'état fondamental au premier état excité.
    \item Représente cette transition à l'aide d'une flèche sur le même diagramme.
    \item Cette énergie est apportée à l'atome par une radiation lumineuse monochromatique.
    \begin{enumerate}[label=c.\arabic*.]
      \item Calcule sa longueur d'onde.
      \item À quel domaine du spectre appartient cette radiation ?
    \end{enumerate}
  \end{enumerate}
  \item Calcule la longueur d'onde de la radiation susceptible d'ioniser l'atome d'hydrogène.
\end{enumerate}
\noindent Données : $c = 3\cdot10^8$ m$\cdot$s$^{-1}$ ; $h = 6{,}62\cdot10^{-34}$ J$\cdot$s ; $1$ eV $= 1{,}6\cdot10^{-19}$ J ; $1$ nm $= 10^{-9}$ m.

\begin{center}
\begin{tikzpicture}[>=Stealth,scale=1]
  \draw[->] (0,0)--(7,0) node[right]{$\lambda$(nm)};
  \draw (2.5,0.08)--(2.5,-0.08) node[below]{400};
  \draw (4.5,0.08)--(4.5,-0.08) node[below]{800};
  \node at (1.2,0.3){\small lumière UV};
  \node at (3.5,0.3){\small lumière visible};
  \node at (5.7,0.3){\small lumière IR};
\end{tikzpicture}
\end{center}

% ============================ PHYSIQUE ============================
\matiere{PHYSIQUE}{12 points}

\partie{A : vérification des connaissances}

\noindent\textbf{1. Schéma à compléter.} Dans le repère ci-contre, représente les vecteurs de Fresnel $\vv{OA_1}$ et $\vv{OA_2}$ associés respectivement aux fonctions sinusoïdales : $y_1 = 2\sin(100\pi t + \pi)$ (cm) et $y_2 = 3\sin\!\left(100\pi t - \dfrac{\pi}{2}\right)$ (cm).

\begin{center}
\begin{tikzpicture}[>=Stealth,scale=1]
  \draw[->] (-2.6,0)--(2.6,0) node[right]{$x$(cm)};
  \draw[->] (0,-2.6)--(0,2.6) node[above]{$y$(cm)};
  \node[above left] at (-2.6,0){$x'$};
  \node[below right] at (0,-2.6){$y'$};
  \draw[->] (0,0)--(0.7,0) node[below]{$\vi$};
  \draw[->] (0,0)--(0,0.7) node[left]{$\vj$};
  \node[below right] at (0,0){$O$};
\end{tikzpicture}
\end{center}

\medskip
\noindent\textbf{2. Question à réponse construite.} Dans l'expérience des fentes de Young en lumière monochromatique, qu'observe-t-on sur l'écran dans la zone où se croisent les lumières issues des sources $S_1$ et $S_2$ ?

\medskip
\noindent\textbf{3. Question à choix multiple.} Choisis la bonne réponse. Exemple : $C = C_2$.
\begin{enumerate}[label=\alph*.]
  \item L'équation différentielle d'un pendule pesant dans le cas des oscillations de faible amplitude s'écrit :
  \begin{enumerate}[label=a.\arabic*.]
    \item $\ddot{\theta} + \dfrac{mg\,OG}{J_\Delta}\theta = 0$ ;
    \item $\ddot{\theta} + \dfrac{J_\Delta}{mg\,OG}\theta = 0$ ;
    \item $\ddot{\theta} + \dfrac{mg\,OG}{J_\Delta}\dot{\theta} = 0$.
  \end{enumerate}
  \item Dans un circuit RLC série, à la résonance :
  \begin{enumerate}[label=b.\arabic*.]
    \item $L\omega$ et $Z$ sont égaux ;
    \item l'intensité efficace est maximale ;
    \item la puissance consommée est minimale.
  \end{enumerate}
\end{enumerate}

\partie{B : application des connaissances}
Un oscilloscope bicourbe permet de visualiser l'intensité dans un circuit RLC et la tension aux bornes du même circuit. On obtient les courbes suivantes : 1 carreau pour 100 V ; 1 carreau pour 100 mA ; 1 carreau pour 1 ms.

\begin{center}
\begin{tikzpicture}[>=Stealth,scale=0.95]
  % grille
  \draw[very thin,gray!40] (0,-2.2) grid[xstep=0.5,ystep=0.5] (6.5,2.2);
  % axes
  \draw[->] (0,0)--(6.7,0) node[right]{$t$(ms)};
  \draw[->] (0,-2.3)--(0,2.4);
  \foreach \k in {1,...,12}{ \draw (\k*0.5,0.06)--(\k*0.5,-0.06) node[below=1pt]{\scriptsize\k}; }
  % u(t) solide : periode ~8ms (4 carreaux)
  \draw[thick] plot[domain=0:6.3,samples=160] (\x,{1.9*sin(deg(2*pi*\x/4))});
  % i(t) tirets : en avance de phase
  \draw[thick,dashed] plot[domain=0:6.3,samples=160] (\x,{1.5*sin(deg(2*pi*\x/4 + 70))});
  \node at (1,2.1){$u(t)$};
  \node[right] at (6.3,1.2){$i(t)$};
\end{tikzpicture}
\end{center}

\begin{enumerate}
  \item \begin{enumerate}[label=\alph*.]
    \item Déduis les valeurs maximales de la tension et de l'intensité.
    \item Calcule l'impédance $Z$ du circuit.
  \end{enumerate}
  \item \begin{enumerate}[label=\alph*.]
    \item Laquelle des deux fonctions $i(t)$ et $u(t)$ est-elle en avance de phase sur l'autre ?
    \item Calcule le déphasage entre les deux fonctions.
  \end{enumerate}
  \item On donne $i(t) = I_m\sin\omega t$. Donne l'expression de la tension $u(t)$ aux bornes du circuit.
\end{enumerate}

\partie{C : résolution d'un problème}
Une particule $\alpha$ (noyau d'hélium $\ch{He^{2+}}$ de charge $q = +2e$) arrive au point $O$ dans un condensateur plan avec une vitesse initiale $\vec{V_0}$ de direction parallèle aux armatures $(C)$ et $(D)$. Une tension constante $U$ est appliquée entre ces deux armatures longues de $L = 5$ cm et distantes de $d = 4$ cm.

\begin{center}
\begin{tikzpicture}[>=Stealth,scale=1]
  % plaques
  \draw[thick] (0,1.3)--(5,1.3) node[midway,above]{$(C)$};
  \draw[thick] (0,-1.3)--(5,-1.3) node[midway,below]{$(D)$};
  % axes
  \draw[->] (0,0)--(6.2,0) node[right]{$x$};
  \draw[->] (0,-0.1)--(0,2) node[above]{$y$};
  % V0
  \draw[->,thick] (0,0)--(1.1,0) node[above]{$\vec{V_0}$};
  \node[below left] at (0,0){$o$};
  % S
  \fill (5,0.6) circle (1.5pt) node[right]{$S$};
  % L
  \draw[<->,dashed] (0,-1.7)--(5,-1.7) node[midway,below]{$L$};
  % d
  \draw[<->,dashed] (-0.35,-1.3)--(-0.35,1.3) node[midway,left]{$d$};
\end{tikzpicture}
\end{center}

\noindent On se propose de déterminer la valeur de la tension $U$ pour que l'ordonnée du point de sortie $S$ soit $Y_S = 1{,}0$ cm.

\begin{enumerate}
  \item Indique la polarité des plaques pour que la particule soit déviée vers le haut.
  \item Recopie la figure en représentant le champ électrique $\vec{E}$ et la force électrique $\vec{F}$ au point $O$.
  \item Établis les équations horaires du mouvement de la particule dans le repère $(Oxy)$.
  \item Déduis l'équation de la trajectoire de la particule à l'intérieur du condensateur.
  \item \begin{enumerate}[label=\alph*.]
    \item Trouve, à partir de l'équation de la trajectoire, l'expression de $U$ en fonction de la masse $m$ de la particule, $e$, $V_0$, $L$ et $Y_S$.
    \item Calcule la valeur de $U$.
  \end{enumerate}
\end{enumerate}
\noindent On donne : $V_0 = 5\cdot10^5$ m$\cdot$s$^{-1}$ ; $e = 1{,}6\cdot10^{-19}$ C ; $m = 6{,}64\cdot10^{-27}$ kg.

\end{document}
